\section{Offline Data Analysis}
\label{sec:analysis}

I ran all offline analyses in Python (NumPy, SciPy, statsmodels) with a pipeline I wrote for this study. The pipeline reads each trial file, corrects artifacts, computes HRV metrics and \al, and stores the results in a local database \todo{versions of Python and packages}. A browser dashboard shows each trial next to the metrics reported in real time by the streaming application.

\subsection{Preprocessing}
\label{sec:preprocessing}

I took RR intervals in the order received and converted them to milliseconds. I identified artifacts in two steps. First, intervals outside 300--2000~ms, corresponding to heart rates outside about 30--200~beats per minute, were marked as physiologically implausible. Second, each remaining interval was compared with the last accepted interval, and intervals differing from it by more than 20\% were marked as artifacts. This percentage rule catches missed beats, extra beats and most ectopic beats. Marked intervals were replaced by linear interpolation between the neighbouring accepted intervals, which keeps the series continuous and preserves its length, as DFA requires.

I computed the proportion of corrected intervals for every trial and report it in Section~\ref{sec:results}. \al{} is more sensitive to artifacts than amplitude indices, and its bias becomes relevant above a few percent \citep{rogers2022update}. I therefore excluded trials with more than 5\% corrected intervals from the analyses of \al{} \confirm{threshold: 5\% vs 3\%}. A sensitivity analysis repeated the main models with a stricter threshold of 3\%.

\subsection{Heart Rate and Amplitude Indices}
\label{sec:hrv-metrics}

For each trial the pipeline computed mean heart rate as $60\,000/\overline{RR}$ (beats per minute), SDNN as the standard deviation of the corrected RR intervals, RMSSD as
\begin{equation}
\mathrm{RMSSD} = \sqrt{\frac{1}{N-1}\sum_{i=1}^{N-1} \left( RR_{i+1} - RR_i \right)^2},
\end{equation}
and pNN50 as the percentage of successive differences larger than 50~ms \citep{taskforce1996}. These indices are reported alongside \al{} to allow comparison with earlier work on flow and skilled performance, which has relied on them, and mean heart rate enters the models as the covariate that indexes physical load.

\subsection{DFA $\alpha_1$}
\label{sec:dfa-implementation}

I computed \al{} following \citet{peng1995}, as described in Section~\ref{sec:bg-dfa}. The corrected RR series was integrated after subtraction of its mean and divided into non-overlapping boxes for each box size $n = 4, 5, \dots, 16$ beats. In each box a linear trend was fitted by least squares and removed, and the fluctuation function $F(n)$ was computed as the root mean square of the residuals across all boxes. \al{} was the slope of a least-squares line fitted to $\log F(n)$ against $\log n$. Estimates were not computed for series shorter than 64 beats, which ensures at least four boxes at the largest box size.

The primary outcome was a single \al{} value per trial, computed over the full two-minute trial. To describe how \al{} evolved within longer recordings, the pipeline also computed a sliding-window estimate over windows of 120 beats moved in steps of 30 beats. These within-trial traces are used descriptively only, because a two-minute trial contains essentially one independent estimate, so within-trial dynamics cannot be tested.

Two settings differ from the reference implementation used in most exercise studies and were aligned with it before the final analysis \confirm{implemented and re-run}. First, smoothness-priors detrending with $\lambda = 500$ \citep{tarvainen2002} was applied to the RR series before DFA, to remove slow trends in heart rate that would otherwise inflate the fluctuation function at the larger box sizes. Second, the sliding-window traces were defined in time (two-minute windows, recomputed every 5~s), so that windows contain comparable durations of data across conditions with different heart rates \citep{rogers2022update}. Results with both the original and the aligned settings are reported, so that the effect of these choices is visible \citep{hoos2025}.

\subsection{Pilot Checks}
\label{sec:pilot-checks}

I used the first recordings to check the pipeline end to end. They confirmed that trials started and stopped at the set duration, that RR intervals arrived without gaps over the WebSocket link, and that the dashboard reproduced the files on disk. They also showed that the real-time \al{} reported by the streaming application could differ noticeably from the offline estimate for the same trial; in one baseline recording, for example, the two values were 1.21 and 1.48. Differences of this size are large relative to the shifts expected between conditions, which is why the primary analyses use the offline estimate only and the agreement between the two is analysed formally (Section~\ref{sec:stats}). \confirm{the first recordings used a seated baseline; record a standing baseline for these participants or analyse them separately} \todo{summarise artifact rates and heart rates from the first five participants once processed}

\subsection{Statistics}
\label{sec:stats}

Table~\ref{tab:analysis-plan} summarises the analysis plan. All models were linear mixed-effects models with a random intercept per participant, fitted by restricted maximum likelihood; fixed effects were tested with Wald tests and reported with 95\% confidence intervals. Movement trials were averaged across hands for H1 and H3 and entered separately for H2, with hand as a covariate. Model assumptions were checked with residual plots; where residuals departed clearly from normality, results were checked with a non-parametric test on per-participant means.

\begin{table}[ht]
\centering
\caption{Analysis plan. $\Delta\alpha_1$ is the \al{} of a movement trial minus the participant's pre-task baseline value.}
\label{tab:analysis-plan}
\small
\begin{tabularx}{\textwidth}{l l X}
\toprule
Question & Outcome & Model \\
\midrule
H1: task vs baseline & \al & condition (baseline, isolation, spin) + random intercept; paired $t$-test on per-participant means as check \\
Direction of shift & $\Delta\alpha_1$ & sign test on the number of participants whose \al{} rises vs falls; position of task values relative to 0.5, 0.75, 1.0 \\
H2: isolation vs spin & \al & movement + mean HR + hand + random intercept (movement trials only) \\
H3: expertise & $|\Delta\alpha_1|$ & expertise score + $\Delta$HR + movement + random intercept; novices vs experienced described separately \\
Carry-over & \al & post-task vs pre-task baseline, paired \\
Hand & \al & dominant vs non-dominant hand, coded by handedness (exploratory) \\
Agreement & \al & Bland--Altman analysis, real-time vs offline estimate \\
\bottomrule
\end{tabularx}
\end{table}

For H1, the contrast of interest was baseline versus movement, with the two movement levels compared with the baseline separately and jointly. Because no direction was predicted, the direction of the shift was reported for every participant, and the proportion of increases and decreases was tested with a sign test. For H2, mean heart rate as a covariate tests whether isolations and spins still differ once their difference in physical load is accounted for. I repeated this analysis with the change in heart rate from baseline as the covariate. For H3, I used the absolute change in \al{} because no direction was predicted. The heart-rate change entered as a covariate, so that a larger change in less experienced participants could not simply reflect greater effort.

I corrected H2 and H3 for multiple comparisons with the Holm procedure. Exploratory analyses (carry-over, hand, agreement, and relations with questionnaire ratings) were not corrected and are reported as estimates with 95\% confidence intervals. I assessed agreement between real-time and offline \al{} with Bland--Altman analysis, reporting the mean difference (bias) and 95\% limits of agreement. I also checked whether the difference depended on the size of \al, as reported for the H10 against ECG \citep{schaffarczyk2022}.
